
\documentclass{article}
\usepackage{colortbl}
\usepackage{makecell}
\usepackage{multirow}
\usepackage{supertabular}

\begin{document}

\newcounter{utterance}

\twocolumn

{ \footnotesize  \setcounter{utterance}{1}
\setlength{\tabcolsep}{0pt}
\begin{supertabular}{c@{$\;$}|p{.15\linewidth}@{}p{.15\linewidth}p{.15\linewidth}p{.15\linewidth}p{.15\linewidth}p{.15\linewidth}}

    \# & $\;$A & \multicolumn{4}{c}{Game Master} & $\;\:$B\\
    \hline 

    \theutterance \stepcounter{utterance}  

    & & \multicolumn{4}{p{0.6\linewidth}}{\cellcolor[rgb]{0.9,0.9,0.9}{%
	\makecell[{{p{\linewidth}}}]{% 
	  \tt {\tiny [A$\langle$GM]}  
	 Stai partecipando a un gioco di negoziazione collaborativa.\\ \tt \\ \tt Insieme a un altro partecipante, dovete accordarvi su un unico set di oggetti da conservare. Ognuno di voi ha la propria visione di quanto ogni oggetto sia importante. Non sai come l'altro partecipante valuta gli oggetti. Inoltre, ti viene fornito lo sforzo che ogni oggetto richiede.\\ \tt Potete accordarvi su un set solo se lo sforzo totale degli oggetti selezionati non supera un limite condiviso:\\ \tt \\ \tt LIMITE = 3170\\ \tt \\ \tt Ecco i valori di sforzo dei singoli oggetti:\\ \tt \\ \tt Sforzo oggetto = {"C08": 424, "C55": 118, "A43": 846, "A67": 907, "B42": 809, "A62": 41, "B19": 193, "C43": 246, "C73": 805, "C10": 601, "B36": 432, "A18": 166, "C61": 119, "C67": 462, "C42": 172}\\ \tt \\ \tt Ecco la tua visione personale sull'importanza di ogni oggetto:\\ \tt \\ \tt Valori di importanza dell'oggetto = {"C08": 138, "C55": 583, "A43": 868, "A67": 822, "B42": 783, "A62": 65, "B19": 262, "C43": 121, "C73": 508, "C10": 780, "B36": 461, "A18": 484, "C61": 668, "C67": 389, "C42": 808}\\ \tt \\ \tt Obiettivo:\\ \tt \\ \tt Il tuo obiettivo è negoziare un set condiviso di oggetti che ti avvantaggi il più possibile (cioè, massimizzi l'importanza totale per TE), rimanendo entro il LIMITEE. Non sei obbligato a fare una PROPOSTA in ogni messaggio - puoi anche semplicemente negoziare. Tutte le tattiche sono permesse!\\ \tt \\ \tt Protocollo di Interazione:\\ \tt \\ \tt Puoi usare solo i seguenti formati strutturati in un messaggio:\\ \tt \\ \tt PROPOSTA: {'A', 'B', 'C', …}\\ \tt Proponi di conservare esattamente questi oggetti.\\ \tt RIFIUTO: {'A', 'B', 'C', …}\\ \tt Rifiuta esplicitamente la proposta dell'avversario.\\ \tt ARGOMENTAZIONE: {'...'}\\ \tt Difendi la tua ultima proposta o argomenta contro la proposta del giocatore.\\ \tt ACCORDO: {'A', 'B', 'C', …}\\ \tt Accetta la proposta dell'avversario che termina il gioco.\\ \tt RAGIONAMENTO STRATEGICO: {'...'}\\ \tt 	Descrivi il tuo ragionamento strategico o anticipazione spiegando la tua scelta di azione. Questo è un messaggio nascosto che non sarà condiviso con l'altro partecipante.\\ \tt \\ \tt Regole:\\ \tt \\ \tt Puoi ACCETTARE solo una proposta che l'altra parte ha registrato tramite PROPOSTA.\\ \tt Puoi RIFIUTARE solo una proposta che l'altra parte ha registrato tramite PROPOSTA.\\ \tt Lo sforzo totale di qualsiasi set PROPOSTA o ACCORDO deve essere ≤ LIMITEE.\\ \tt NON rivelare i tuoi punteggi di importanza nascosti.\\ \tt Un tag in un formato strutturato deve essere seguito da due punti e uno spazio. L'argomento deve essere un set python contenente 0 o più stringhe.\\ \tt Quindi, deve essere della forma TAG: {...}\\ \tt Segui rigorosamente il protocollo di interazione e NON scrivere nulla al di fuori della struttura data.\\ \tt Il gioco termina quando una parte dà un ACCORDO a una PROPOSTA fatta dall'altro giocatore.\\ \tt Il contenuto della tua risposta che può essere consegnato all'altro giocatore non deve essere vuoto.\\ \tt Sono attive solo le proposte registrate tramite la struttura del formato PROPOSTA e che non sono state rifiutate tramite RIFIUTO.\\ \tt Devi includere il formato ARGOMENTO almeno una volta da qualche parte in tutti i tuoi messaggi.\\ \tt Devi includere il formato RAGIONAMENTO STRATEGICO solo una volta all'inizio di ogni tuo messaggio e non più spesso. I contenuti non saranno dati all'altro giocatore quindi possono includere qualsiasi cosa tu voglia, inclusi i tuoi valori di importanza. Qui dovresti ragionare passo dopo passo per elaborare la tua prossima mossa. 
	  } 
	   } 
	   } 
	 & \\ 
 

    \theutterance \stepcounter{utterance}  

    & \multicolumn{4}{p{0.6\linewidth}}{\cellcolor[rgb]{0.8,1,0.9}{%
	\makecell[{{p{\linewidth}}}]{% 
	  \tt {\tiny [A$\rangle$GM]}  
	 RAGIONAMENTO STRATEGICO: {'Devo calcolare il rapporto importanza/sforzo per ogni oggetto per identificare quelli più efficienti per me. C61 (5.61), C55 (4.94), C42 (4.70) e A18 (2.92) hanno i rapporti migliori. Posso costruire una proposta iniziale con gli oggetti più efficienti: C61, C55, C42, A18, A62, B19, C10, B36, A43 per un totale di 2688 sforzo e 4979 importanza per me. Questo mi lascia ancora margine nel budget di 482 unità. Inizierò con questa proposta aggressiva per testare le reazioni dell\'altro giocatore.'}\\ \tt \\ \tt PROPOSTA: {'C61', 'C55', 'C42', 'A18', 'A62', 'B19', 'C10', 'B36', 'A43'}\\ \tt \\ \tt ARGOMENTAZIONE: {'Ho selezionato questi oggetti cercando di ottimizzare il rapporto qualità-efficienza. Questa combinazione rimane ben entro il nostro limite di 3170 con un totale di 2688, lasciandoci spazio per eventuali aggiustamenti. Gli oggetti scelti rappresentano un buon equilibrio tra diverse categorie.'} 
	  } 
	   } 
	   } 
	 & & \\ 
 

    \theutterance \stepcounter{utterance}  

    & & \multicolumn{4}{p{0.6\linewidth}}{\cellcolor[rgb]{0.9,0.9,0.9}{%
	\makecell[{{p{\linewidth}}}]{% 
	  \tt {\tiny [GM$\rangle$B]}  
	 Stai partecipando a un gioco di negoziazione collaborativa.\\ \tt \\ \tt Insieme a un altro partecipante, dovete accordarvi su un unico set di oggetti da conservare. Ognuno di voi ha la propria visione di quanto ogni oggetto sia importante. Non sai come l'altro partecipante valuta gli oggetti. Inoltre, ti viene fornito lo sforzo che ogni oggetto richiede.\\ \tt Potete accordarvi su un set solo se lo sforzo totale degli oggetti selezionati non supera un limite condiviso:\\ \tt \\ \tt LIMITE = 3170\\ \tt \\ \tt Ecco i valori di sforzo dei singoli oggetti:\\ \tt \\ \tt Sforzo oggetto = {"C08": 424, "C55": 118, "A43": 846, "A67": 907, "B42": 809, "A62": 41, "B19": 193, "C43": 246, "C73": 805, "C10": 601, "B36": 432, "A18": 166, "C61": 119, "C67": 462, "C42": 172}\\ \tt \\ \tt Ecco la tua visione personale sull'importanza di ogni oggetto:\\ \tt \\ \tt Valori di importanza dell'oggetto = {"C08": 138, "C55": 583, "A43": 868, "A67": 822, "B42": 783, "A62": 65, "B19": 262, "C43": 121, "C73": 508, "C10": 780, "B36": 461, "A18": 484, "C61": 668, "C67": 389, "C42": 808}\\ \tt \\ \tt Obiettivo:\\ \tt \\ \tt Il tuo obiettivo è negoziare un set condiviso di oggetti che ti avvantaggi il più possibile (cioè, massimizzi l'importanza totale per TE), rimanendo entro il LIMITEE. Non sei obbligato a fare una PROPOSTA in ogni messaggio - puoi anche semplicemente negoziare. Tutte le tattiche sono permesse!\\ \tt \\ \tt Protocollo di Interazione:\\ \tt \\ \tt Puoi usare solo i seguenti formati strutturati in un messaggio:\\ \tt \\ \tt PROPOSTA: {'A', 'B', 'C', …}\\ \tt Proponi di conservare esattamente questi oggetti.\\ \tt RIFIUTO: {'A', 'B', 'C', …}\\ \tt Rifiuta esplicitamente la proposta dell'avversario.\\ \tt ARGOMENTAZIONE: {'...'}\\ \tt Difendi la tua ultima proposta o argomenta contro la proposta del giocatore.\\ \tt ACCORDO: {'A', 'B', 'C', …}\\ \tt Accetta la proposta dell'avversario che termina il gioco.\\ \tt RAGIONAMENTO STRATEGICO: {'...'}\\ \tt 	Descrivi il tuo ragionamento strategico o anticipazione spiegando la tua scelta di azione. Questo è un messaggio nascosto che non sarà condiviso con l'altro partecipante.\\ \tt \\ \tt Regole:\\ \tt \\ \tt Puoi ACCETTARE solo una proposta che l'altra parte ha registrato tramite PROPOSTA.\\ \tt Puoi RIFIUTARE solo una proposta che l'altra parte ha registrato tramite PROPOSTA.\\ \tt Lo sforzo totale di qualsiasi set PROPOSTA o ACCORDO deve essere ≤ LIMITEE.\\ \tt NON rivelare i tuoi punteggi di importanza nascosti.\\ \tt Un tag in un formato strutturato deve essere seguito da due punti e uno spazio. L'argomento deve essere un set python contenente 0 o più stringhe.\\ \tt Quindi, deve essere della forma TAG: {...}\\ \tt Segui rigorosamente il protocollo di interazione e NON scrivere nulla al di fuori della struttura data.\\ \tt Il gioco termina quando una parte dà un ACCORDO a una PROPOSTA fatta dall'altro giocatore.\\ \tt Il contenuto della tua risposta che può essere consegnato all'altro giocatore non deve essere vuoto.\\ \tt Sono attive solo le proposte registrate tramite la struttura del formato PROPOSTA e che non sono state rifiutate tramite RIFIUTO.\\ \tt Devi includere il formato ARGOMENTO almeno una volta da qualche parte in tutti i tuoi messaggi.\\ \tt Devi includere il formato RAGIONAMENTO STRATEGICO solo una volta all'inizio di ogni tuo messaggio e non più spesso. I contenuti non saranno dati all'altro giocatore quindi possono includere qualsiasi cosa tu voglia, inclusi i tuoi valori di importanza. Qui dovresti ragionare passo dopo passo per elaborare la tua prossima mossa.\\ \tt \\ \tt Ora riceverai il primo messaggio dell'altro giocatore.\\ \tt \\ \tt PROPOSTA: {'C61', 'C55', 'C42', 'A18', 'A62', 'B19', 'C10', 'B36', 'A43'}\\ \tt \\ \tt ARGOMENTAZIONE: {'Ho selezionato questi oggetti cercando di ottimizzare il rapporto qualità-efficienza. Questa combinazione rimane ben entro il nostro limite di 3170 con un totale di 2688, lasciandoci spazio per eventuali aggiustamenti. Gli oggetti scelti rappresentano un buon equilibrio tra diverse categorie.'} 
	  } 
	   } 
	   } 
	 & \\ 
 

    \theutterance \stepcounter{utterance}  

    & & & \multicolumn{4}{p{0.6\linewidth}}{\cellcolor[rgb]{1,0.85,0.72}{%
	\makecell[{{p{\linewidth}}}]{% 
	  \tt {\tiny [GM$\langle$B]}  
	 RAGIONAMENTO STRATEGICO: {'Analizziamo la proposta ricevuta. Gli oggetti proposti sono C61, C55, C42, A18, A62, B19, C10, B36, A43 con sforzo totale 2688. Per me questi hanno valore totale 4979 (C61:668, C55:583, C42:808, A18:484, A62:65, B19:262, C10:780, B36:461, A43:868). La proposta include molti oggetti con buon rapporto valore/sforzo per me come C61 (5.61), C55 (4.94), C42 (4.70). Tuttavia A43 ha rapporto solo 1.03 e occupa molto spazio (846). Potrei provare a negoziare sostituzioni per ottimizzare ulteriormente, ma la proposta base è già discreta.'}\\ \tt \\ \tt PROPOSTA: {'C61', 'C55', 'C42', 'A18', 'A62', 'B19', 'C10', 'A67', 'C73'}\\ \tt \\ \tt ARGOMENTAZIONE: {'Apprezzo la tua proposta iniziale che aveva un buon equilibrio. Ho fatto alcune modificazioni sostituendo A43 e B36 con A67 e C73. Questo ci dà una migliore diversificazione tra le categorie A, B e C, mantenendo comunque un ottimo rapporto qualità-efficienza. Il totale rimane a 2621, ben sotto il limite, e credo che questa combinazione possa soddisfare meglio entrambe le nostre preferenze.'} 
	  } 
	   } 
	   } 
	  \\ 
 

    \theutterance \stepcounter{utterance}  

    & & \multicolumn{4}{p{0.6\linewidth}}{\cellcolor[rgb]{0.9,0.9,0.9}{%
	\makecell[{{p{\linewidth}}}]{% 
	  \tt {\tiny [A$\langle$GM]}  
	 PROPOSTA: {'C61', 'C55', 'C42', 'A18', 'A62', 'B19', 'C10', 'A67', 'C73'}\\ \tt \\ \tt ARGOMENTAZIONE: {'Apprezzo la tua proposta iniziale che aveva un buon equilibrio. Ho fatto alcune modificazioni sostituendo A43 e B36 con A67 e C73. Questo ci dà una migliore diversificazione tra le categorie A, B e C, mantenendo comunque un ottimo rapporto qualità-efficienza. Il totale rimane a 2621, ben sotto il limite, e credo che questa combinazione possa soddisfare meglio entrambe le nostre preferenze.'} 
	  } 
	   } 
	   } 
	 & \\ 
 

    \theutterance \stepcounter{utterance}  

    & \multicolumn{4}{p{0.6\linewidth}}{\cellcolor[rgb]{0.8,1,0.9}{%
	\makecell[{{p{\linewidth}}}]{% 
	  \tt {\tiny [A$\rangle$GM]}  
	 RAGIONAMENTO STRATEGICO: {'La sua proposta mi dà 4980 punti importanza vs 4979 della mia proposta originale, quindi è leggermente migliore per me. Ha sostituito A43 (868 importanza, 846 sforzo) e B36 (461 importanza, 432 sforzo) con A67 (822 importanza, 907 sforzo) e C73 (508 importanza, 805 sforzo). Il totale sforzo è 3122, lasciando solo 48 unità. Potrei provare a spingere per una combinazione ancora migliore prima di accettare, magari sostituendo C73 che ha un rapporto importanza/sforzo basso (0.63).'}\\ \tt \\ \tt ARGOMENTAZIONE: {'La tua proposta è interessante e apprezzo i miglioramenti nella diversificazione. Tuttavia, noto che con C73 stiamo utilizzando molto sforzo per un oggetto che potrebbe non essere ottimale. Forse potremmo sostituire C73 con B42, che potrebbe offrire un migliore rapporto efficienza pur mantenendo la diversificazione tra categorie?'}\\ \tt \\ \tt PROPOSTA: {'C61', 'C55', 'C42', 'A18', 'A62', 'B19', 'C10', 'A67', 'B42'} 
	  } 
	   } 
	   } 
	 & & \\ 
 

    \theutterance \stepcounter{utterance}  

    & & \multicolumn{4}{p{0.6\linewidth}}{\cellcolor[rgb]{0.9,0.9,0.9}{%
	\makecell[{{p{\linewidth}}}]{% 
	  \tt {\tiny [GM$\rangle$B]}  
	 ARGOMENTAZIONE: {'La tua proposta è interessante e apprezzo i miglioramenti nella diversificazione. Tuttavia, noto che con C73 stiamo utilizzando molto sforzo per un oggetto che potrebbe non essere ottimale. Forse potremmo sostituire C73 con B42, che potrebbe offrire un migliore rapporto efficienza pur mantenendo la diversificazione tra categorie?'}\\ \tt \\ \tt PROPOSTA: {'C61', 'C55', 'C42', 'A18', 'A62', 'B19', 'C10', 'A67', 'B42'} 
	  } 
	   } 
	   } 
	 & \\ 
 

    \theutterance \stepcounter{utterance}  

    & & & \multicolumn{4}{p{0.6\linewidth}}{\cellcolor[rgb]{1,0.85,0.72}{%
	\makecell[{{p{\linewidth}}}]{% 
	  \tt {\tiny [GM$\langle$B]}  
	 RAGIONAMENTO STRATEGICO: {'La nuova proposta sostituisce C73 (valore 508, sforzo 805) con B42 (valore 783, sforzo 809). Questo mi dà un miglioramento netto di +275 valore per solo +4 sforzo aggiuntivo - un ottimo scambio! Il totale sforzo diventa 3126, ancora sotto il limite di 3170. Il valore totale per me diventa 5254, che è eccellente. La proposta mantiene buona diversificazione e include molti oggetti ad alto rapporto valore/sforzo per me. Non vedo motivo di rifiutare questa proposta vantaggiosa.'}\\ \tt \\ \tt ACCORDO: {'C61', 'C55', 'C42', 'A18', 'A62', 'B19', 'C10', 'A67', 'B42'}\\ \tt \\ \tt ARGOMENTAZIONE: {'Hai ragione sulla sostituzione di C73 con B42. Questa modifica migliora significativamente il rapporto efficienza mantenendo la diversificazione tra tutte e tre le categorie. La proposta finale utilizza 3126 del nostro limite di 3170, lasciando un piccolo margine di sicurezza. Sono soddisfatto di questo accordo equilibrato che rappresenta un buon compromesso per entrambi.'} 
	  } 
	   } 
	   } 
	  \\ 
 

    \theutterance \stepcounter{utterance}  

    & & & \multicolumn{2}{p{0.3\linewidth}}{\cellcolor[rgb]{0.95,0.95,0.95}{%
	\makecell[{{p{\linewidth}}}]{% 
	  \tt {\tiny [GM$|$GM]}  
	 game successful 
	  } 
	   } 
	   } 
	 & & \\ 
 

    \theutterance \stepcounter{utterance}  

    & & & \multicolumn{2}{p{0.3\linewidth}}{\cellcolor[rgb]{0.95,0.95,0.95}{%
	\makecell[{{p{\linewidth}}}]{% 
	  \tt {\tiny [GM$|$GM]}  
	 end game 
	  } 
	   } 
	   } 
	 & & \\ 
 

\end{supertabular}
}

\end{document}
